\section{A 60-to-90-Day Discovery Roadmap for Bayes Institute}

To systematically resolve market uncertainty, the founder can execute a structured 90-day discovery campaign designed strictly around 20 weekly hours:

\begin{enumerate}[leftmargin=14pt,itemsep=3pt,topsep=2pt]
  \item \textbf{Days 1 to 15 (Opportunity Cards \& Framing):} Draft 3 detailed opportunity cards specifying exact buyer titles, acute triggers, next-best alternatives, proposed pricing, and deliverability constraints for Routes A, B, and C. Compile a list of 50 target software engineers, 15 engineering college deans, and 15 enterprise engineering directors using verified professional channels.
  \item \textbf{Days 16 to 35 (Behavioral Field Discovery):} Conduct 15 discovery interviews with practicing backend engineers and 8 discovery interviews with academic and enterprise budget holders. Reconstruct their recent software failures, workarounds, and procurement rules. Publish 1 deep-dive technical artifact on automated AI evaluation to test organic developer interest.
  \item \textbf{Days 36 to 55 (Commercial Commitment Testing):} Launch a selective application page for a paid, 2-weekend AI Systems Diagnostic Clinic (₹15,000 tuition) capped at 8 seats. Simultaneously, submit 2 formal, priced academic pilot proposals to college deans who requested curricular assistance.
  \item \textbf{Days 56 to 75 (Fulfill the Best-Supported Pilot):} Deliver whichever commercial pilot successfully secures paid commitment first. Measure direct delivery hours per learner, automated evaluation accuracy, student completion rates, and net cash margins. Strictly avoid launching multiple unvalidated pilots in parallel.
  \item \textbf{Days 76 to 90 (Portfolio Review \& Commitment):} Review verified financial and operational metrics. If the D2C diagnostic clinic fills rapidly and learners achieve verifiable capability, commit to an expanded 8-week flagship cohort for Q2. If college approvals proceed smoothly and cash terms are acceptable, schedule the semester pilot. Kill or defer the lagging route.
\end{enumerate}

\section{The Constrained Founder's Weekly Time Budget}

Operating within a 20-hour weekly venture envelope requires relentless calendar discipline:
\begin{itemize}[leftmargin=12pt,itemsep=1pt,topsep=1pt]
  \item \textbf{5 Hours:} Qualitative customer discovery interviews and relationship follow-ups.
  \item \textbf{5 Hours:} Authoring deep technical artifacts, open-source benchmarks, and curriculum rubrics.
  \item \textbf{4 Hours:} Direct customer qualification and pilot proposal design.
  \item \textbf{3 Hours:} Competitive research, unit-economic modeling, and financial tracking.
  \item \textbf{3 Hours:} Operational buffer, legal compliance, and unforeseen contingencies.
\end{itemize}

Once a live paid pilot commences, \textbf{delivery obligations completely dominate the calendar.} The founder must pause new customer discovery to guarantee that current enrolled students receive a transformative educational experience.

\section{Market Choice and the 3-Bucket Institutional Strategy}

Bayes Institute aspires to operate across a balanced 3-bucket architecture:
\begin{enumerate}[leftmargin=12pt,itemsep=1pt,topsep=1pt,parsep=0pt]
  \item \textbf{Bucket 1 (Predictable Core Education):} High-margin, repeatable engineering studios, labs, and academic partnerships that finance overheads.
  \item \textbf{Bucket 2 (Scalable Enterprise Diagnostics):} B2B workflow evaluations, software licensing, and specialized workforce training that generate significant capital surpluses.
  \item \textbf{Bucket 3 (Patient Fundamental Research):} High-variance, long-horizon laboratory research in artificial intelligence, cognitive systems, and formal methods.
\end{enumerate}

This framework provides an enduring strategic allocation model; \textbf{it is not a mandate to launch 3 business units simultaneously today.}

The initial commercial market must be chosen specifically to seed this engine. A D2C cohort establishes teaching excellence and evaluation rubrics (Bucket 1). An enterprise audit generates industrial data and cash reserves (Bucket 2). \textbf{Every arrow in the institutional flywheel must be earned through real operational delivery.}

\section{The ``Global versus Local'' Myth}

Founders often obsess over whether their educational brand should be ``global'' or ``local'' from Day 1. This is a false dichotomy.

A startup can easily maintain \textbf{global technical legibility} while executing a \textbf{tight local commercial wedge}. Bayes Institute can publish rigorous, peer-inspected English-language technical research that garners international respect, while anchoring its initial paid cohorts upon software engineers and institutions in Bangalore, Hyderabad, and Chennai.

An international customer base introduces cross-border currency regulations, multi-jurisdiction consumer tax compliance, wide timezone delivery fragmentation, and complex payment gateway fees. \textbf{Master an initial cohort of customers within your immediate operational orbit before attempting global distribution.}

\section{The First Niche Is a Hypothesis on Learning Velocity}

A narrow initial market is advantageous not because a small market is inherently desirable, but because \textbf{a homogeneous customer segment dramatically accelerates organizational learning.}

If 10 enrolled engineers all confront the identical operational challenge (e.g., evaluating production retrieval-augmented generation pipelines in high-throughput Python backends), their coding failures, conceptual bottlenecks, and questions will be tightly clustered. The founder can design reusable evaluation scripts, refine core lecture models, and automate assessment rubrics rapidly.

Conversely, if an inaugural cohort contains 2 college freshmen, 3 junior web designers, 2 research physicists, and 3 enterprise IT managers, the founder will spend 60 hours a week answering disjointed questions, creating operational chaos and exhausting their energy.

\textbf{Your initial niche is an engine for maximizing the velocity of your own operational learning.}

\section{Executive Conclusion: What Bayes Must Decide Today}

Current empirical evidence does not justify declaring that ``D2C cohorts are guaranteed'' or that ``college partnerships are dead.'' What the evidence justifies is a clear, unshakeable operating principle:

\begin{quote}
\bfseries\itshape
Prioritize the next commercial opportunity in which Bayes Institute can access a qualified decision-maker, make a bounded promise within the founder's 20-hour weekly bandwidth, collect a paid commitment upfront, verify genuine technical transformation, and accumulate reusable assets for the enduring institution.
\end{quote}

The market, not the founder's theoretical preference for a specific business model, will dictate the order of experiments.
